\section{Related Work}

\smallsec{Theorem Proving}
Classical provers express theorems in first-order logic and search for proofs automatically in a large space~\cite{robinson2001handbook,kovacs2013first}. Even with data-driven search heuristics~\cite{loos2017deep,bridge2014machine}, they fail to scale to large formalization projects. Therefore, recent work on learning-based theorem proving has focused on an alternative paradigm: automating the interaction with proof assistants. 

The architecture of learning-based provers progressed from classical machine learning algorithms such as KNN~\cite{gauthier2021tactictoe}, to graph neural networks explicitly encoding the syntax of formal expressions~\cite{yang2019learning,paliwal2020graph}, and now Transformer-based LLMs treating expressions as plain strings~\cite{polu2020generative}. Besides the model architecture, researchers have explored several complementary dimensions: proof search algorithms for assembling model-generated steps into complete proofs~\cite{lample2022hypertree,wang2023dt}; overcoming data scarcity through reinforcement learning (RL)~\cite{bansal2019learning,wu2021tacticzero,lample2022hypertree,polu2023formal} or synthetic/auxiliary data~\cite{wang2020learning,rabe2021mathematical,wu2021lime,han2022proof}; as well as outsourcing some proof goals to classical provers~\cite{bohme2010sledgehammer,blanchette2016hammering,czajka2018hammer,jiang2022thor}. Our base model without retrieval is a combination of straightforward design choices. It generates tactics by finetuning an encoder-decoder Transformer, ByT5~\cite{xue2022byt5}, via supervised learning without RL or auxiliary data. Then it searches for proofs using best-first search. Our model's algorithmic novelty lies in the retrieval.


\smallsec{Premise Selection}
Selecting useful premises is recognized as a key challenge in theorem proving~\cite{urban2004mptp,irving2016deepmath,szegedy2021retrieval,wu2022formal}. Machine learning methods for premise selection have also progressed from classical models~\cite{alama2014premise,bohme2010sledgehammer,piotrowski2023machinelearned}, recurrent neural networks~\cite{irving2016deepmath}, graph neural networks~\cite{wang2020learning}, to Transformers~\cite{mikula2023magnushammer,yeh2023coprover}. However, existing methods either tackle premise selection in isolation without theorem proving~\cite{piotrowski2023machinelearned,irving2016deepmath,wang2020learning} or feed the premises to a symbolic prover~\cite{alama2014premise,bohme2010sledgehammer,mikula2023magnushammer}. To our knowledge, we are the first to augment a learning-based formal theorem prover with retrieved premises so that the prover can learn how to use them effectively. For example, it can decide whether to use an explicitly retrieved premise or an implicitly memorized one.


\smallsec{Data and Tools for Theorem Proving}
Tools for data extraction and interacting with proof assistants have been crucial drivers of learning-based theorem proving. Existing tools and datasets can be divided by proof assistants: Coq has GamePad~\cite{huanggamepad}, CoqGym~\cite{yang2019learning}, and PRISM~\cite{reichel2023proof}; Isabelle has IsarStep~\cite{li2021isarstep} and PISA~\cite{jiang2021lisa}; HOL Light has HOList~\cite{bansal2019holist} and HoLStep~\cite{kaliszykholstep}, and Lean has LeanStep~\cite{han2022proof} and \texttt{lean-gym}~\cite{polu2023formal}. MiniF2F~\cite{zheng2022minif2f} is the only cross-system dataset, with 488 theorems for evaluation. However, it does not have training theorems and is restricted to the domain of math olympiads. 
 
Among available tools extracting data from proof assistants, LeanDojo is the only one that can extract premises for retrieval-augmented theorem proving. A few existing datasets also have premises~\cite{bansal2019holist,mikula2023magnushammer}, but their data extraction tools are not public, making it difficult to construct new datasets. In addition, LeanDojo is the only tool that can interact with Lean robustly (Sec.~\ref{sec:leandojo}) and can extract data from Lean 4. See Appendix~\ref{appendix:tool_comparison} for a detailed comparison between LeanDojo and alternatives.


\smallsec{Mathematical Reasoning in Natural Language}
We focus on proving theorems expressed in formal logic, whereas researchers have also produced a plethora of work on mathematical reasoning in natural language~\cite{hendrycks2021measuring,lewkowycz2022solving,ferreira2020premise,welleck2021naturalproofs,welleck2022naturalprover,mishra2022lila,lu2022survey,meadows2022survey}. A particularly relevant task is autoformalization, translating natural language texts into formal theorems and proofs~\cite{wang2018first,cosler2023nl2spec,pan2023data,hahn2022formal,wu2022autoformalization,jiang2023draft,zhao2023decomposing,cunningham2023towards,azerbayev2023proofnet,chen2023nl2tl}. 

\smallsec{Retrieval-Augmented Language Models}
Our {\name} is the first retrieval-augmented language model for formal theorem proving, though similar architectures have been studied extensively in NLP~\cite{khandelwal2020generalization,guu2020retrieval,lewis2020retrieval,borgeaud2022improving,li2022decoupled,jiang2022retrieval,wu2022memorizing,izacard2022few,zhong2022training}. In addition, there have been many retrieval-augmented methods for code generation~\cite{hayati2018retrieval,parvez2021retrieval,lu2022reacc,zhou2023docprompting,shrivastava2022repository,zhang2023repocoder,ding2022cocomic}. Most of them retrieve from a corpus not directly related to the current file, e.g., GitHub or Stack Overflow. In contrast, our retrieval corpus consists of premises accessible to the current file, which is determined by program analysis using LeanDojo. This is similar to what CoCoMIC~\cite{ding2022cocomic} does for Python. However, their retrieval is based on heuristics, whereas ours is learned. 